\section{Context: from 2.0 to 3.0}
\label{sec:context}

\Zoo{} 3.0 is the post-quantum hardening epoch of the \Zoo{} stack.
This section lays out where \Zoo{} stood at the start of 2025, what
\Zoo{} 2.0 left unfinished, and what 3.0 commits to deliver.

\subsection{The arc up to 3.0}

\paragraph{1.0 (2021, BSC).}
The original \Zoo{} whitepaper of October 31, 2021 \cite{zoo-2021}
positioned \Zoo{} as a wildlife-conservation platform on Binance Smart
Chain. That paper coined the term ``NFT Liquidity Protocol'' --- the
mechanism by which \NFT{}s become productive collateral inside a
yield-bearing system. Cryptographic posture was inherited from BSC:
ECDSA over secp256k1, Ethereum-style Merkle Patricia tries, no
post-quantum guarantees.

\paragraph{2.0 (October 2024, Lux L2).}
\Zoo{} 2.0 \cite{zoo-2-0} re-launched on the \Lux{} Network as an
\textsc{evm}-compatible L2. The 2.0 paper introduced 20 native
precompiles including NIST PQC standards ($\MLDSA{}$, SLH-DSA,
ML-KEM), threshold signatures (FROST, CGGMP21), privacy primitives
(\FHE{}, \ZKP{}, ring signatures), and an AI Mining precompile. 2.0
delivered \emph{first-pass} post-quantum: PQ primitives were available
to applications via precompile, but the chain consensus, governance,
bridge, and the legacy 2021 \NFT{} Liquidity Protocol commitments
still relied on classical cryptography for some edges.

\paragraph{3.0 (2025, full PQ).}
\Zoo{} 3.0 closes those edges. The contract: every signature on the
hot path, every commitment in governance, every proof in the bridge,
every receipt under the \NFT{} Liquidity Protocol uses post-quantum
schemes. PQ is no longer optional or per-application; it is the
chain-wide default.

\paragraph{4.0 (2026, GPU-native L1).}
\Zoo{} 4.0 \cite{zoo-4-0} graduates \Zoo{} from L2 application to
sovereign L1 with $100\%$ \GPU{}-powered consensus, \textsc{evm},
\DEX{}, AI mining, and \FHE{}. 4.0 inherits 3.0's full-PQ posture
unchanged.

\subsection{The threat model that drives 3.0}

\Zoo{}'s threat model in 3.0 is the harvest-now-decrypt-later
adversary plus a near-term Q-day adversary. Harvest-now is already
operational: a passive observer can snapshot any classical signature
ever recorded on a public chain and decrypt years later when a
cryptographically relevant quantum computer (CRQC) becomes available.
This means classical signatures recorded today on hot-path data
(governance commitments, bridge attestations, \NFT{} loan agreements)
already need to be PQ-safe, because the value those signatures
authenticate may persist far past Q-day.

A near-term Q-day adversary further means the chain must be able
to validate the integrity of pre-Q-day records using post-Q-day
verifiers. \Zoo{} 3.0's triple-cert design (see \S\ref{sec:triple-cert})
addresses this directly: every block carries a classical proof and
two independent post-quantum proofs, so even when the classical layer
is broken the post-quantum layers continue to validate.

\subsection{What 3.0 changes (one-page summary)}

\begin{longtable}{p{0.32\textwidth}p{0.30\textwidth}p{0.32\textwidth}}
\toprule
\textbf{Component} & \textbf{2.0 (October 2024)} & \textbf{3.0 (2025)} \\
\midrule
Consensus signatures & inherited L2 from \Lux{} 2.x & triple cert via \Quasar{} 3.0 \cite{lp-020} \\
Governance keys & ECDSA per signer & threshold $\MLDSA{}\text{-}65$ + \Corona{} \\
Bridge attestations & classical multisig & threshold \MLDSA{} on M-Chain (\textsc{lp-016}, \textsc{lp-017}) \\
zLLM commitments & SHA-256 Merkle & PQ-safe accumulator + \MLDSA{}-signed root \\
NFT Liquidity loans & ECDSA loan agreements & \Corona{}-signed loan agreements \\
DAO vote receipts & on-chain ECDSA & \MLDSA{}-65 signed, optional \FHE{} ballot \\
Inference receipts & A-Chain TEE quote & TEE quote + \Corona{}-anchored attribution \\
Wallet signing keys & secp256k1 only & dual-key (secp256k1 + \MLDSA{}-65) \\
\bottomrule
\end{longtable}

The remaining sections expand each row into its specification and
migration plan.

\subsection{What 3.0 explicitly does \emph{not} change}

\begin{itemize}[leftmargin=1.2em]
  \item \Zoo{} stays an L2 on \Lux{}; the L1 graduation is a 4.0
        concern \cite{zoo-4-0}.
  \item \textsc{evm} bytecode semantics are unchanged; only the
        signature scheme over transactions and state commitments
        moves to PQ-safe forms.
  \item Application contracts written for 2.0 continue to work in 3.0;
        new precompiles for \Corona{} verification are additive,
        not breaking.
  \item Token standards (ZRC-20, ZRC-721, ZRC-1155) are unchanged; only
        the underlying signing key types are extended.
\end{itemize}

\subsection{Why 2025 is the right epoch for full PQ}

NIST finalised FIPS 204 ($\MLDSA{}$) in August 2024 \cite{nist-fips-204},
which gives \Zoo{} a standardised lattice signature with broad
ecosystem support. The \Lux{} Network shipped \Quasar{} 3.0 triple
consensus \cite{lp-020} in 2025, giving \Zoo{} a chain-level
substrate for triple-cert proofs without per-app integration. The
remaining work in 2025 is the application-layer migration: governance,
bridge, NFT Liquidity, zLLM commitments, wallet keys. \Zoo{} 3.0
sequences that migration through the rest of the year so that 4.0
inherits a uniformly post-quantum stack.
